\par\Needspace{4\baselineskip}
\begin{center}
\hypertarget{native-035-article017}{}\pdfbookmark[2]{17.}{article-017}17.
\end{center}
It is worth adding a few remarks on finding the \emph{maximum equation of centre}. First, the difference between the eccentric and mean anomalies is plainly greatest at $E=90^0$, where it is $=e$ (to be expressed in degrees and their subdivisions). At this point the radius vector is $=a$, so $v=90^0+\varphi$ and the full equation of centre is therefore $=\varphi+e$. It does not yet attain its maximum, since the difference between $v$ and $E$ can still increase beyond $\varphi$. \emph{This} difference is greatest when $\mathrm{d}(v-E)=0$, or when $\mathrm{d}v=\mathrm{d}E$, the eccentricity clearly being held constant. Under this assumption we have $\frac{\mathrm{d}v}{\sin v}=\frac{\mathrm{d}E}{\sin E}$ in general. Thus at the point where the difference between $v$ and $E$ is greatest, we must have $\sin v=\sin E$. Hence, by equations VIII and III, we have
\[
r=a\cos\varphi
\]
\[
e\cos E=1-\cos\varphi
\]
or
\[
\cos E=+\operatorname{tang}\frac{1}{2}\varphi.
\]
Likewise we find $\cos v=-\operatorname{tang}\frac{1}{2}\varphi$, so we have\footnote{The maxima between aphelion and perihelion may be left aside, since they plainly differ from those between perihelion and aphelion only in sign.}
\[
v=90^0+\operatorname{arc.}\sin\operatorname{tang}\frac{1}{2}\varphi
\]
\[
E=90^0-\operatorname{arc.}\sin\operatorname{tang}\frac{1}{2}\varphi;
\]
It follows further that $\sin E=\BareRoot{(1-\operatorname{tang}\frac{1}{2}\varphi^2)}=\frac{\BareRoot{\cos\varphi}}{\cos\frac{1}{2}\varphi}$, so the full equation of centre at this point becomes
\[
=2\operatorname{arc.}\sin\operatorname{tang}\frac{1}{2}\varphi+2\sin\frac{1}{2}\varphi\cdot\BareRoot{\cos\varphi}
\]
with the second term expressed in degrees and their subdivisions.

Finally, at the point where the full equation of centre itself is greatest, we must have $\mathrm{d}v=\mathrm{d}M$ and hence, by Article 15, $r=a\BareRoot{\cos\varphi}$. It follows that
\SourceLeaf{036}
\[
\cos v=-\frac{1-\cos\varphi^{\frac{3}{2}}}{e}
\]
\[
\cos E=\frac{1-\BareRoot{\cos\varphi}}{e}=\frac{1-\cos\varphi}{e(1+\BareRoot{\cos\varphi})}=\frac{\operatorname{tang}\frac{1}{2}\varphi}{1+\BareRoot{\cos\varphi}},
\]
from which $E$ can be determined with the highest precision. Once $E$ is known, the equation of centre is, by equations X and XII, $=2\operatorname{arc.}\sin\frac{\sin\frac{1}{2}\varphi\sin E}{\BareRoot[4]{\cos\varphi}}+e\sin E$\HistoricalNoteEight{Manuscript note:] $\sin E=\BareRoot{\frac{\cos\varphi+2\BareRoot{\cos\varphi}+\cos\varphi^{\frac{3}{2}}}{(1+\BareRoot{\cos\varphi})(1+\cos\varphi)}},$ Equation of centre $=2\operatorname{arc.}\sin\BareRoot{\left(\frac{1}{1+\cos\varphi}-\frac{\BareRoot[4]{(1-e e)}}{2}\right)}+e\sin E=\operatorname{arc.}\cos\frac{-1+\BareRoot{\cos\varphi}+\cos\varphi+\cos\varphi^{\frac{3}{2}}}{1+\cos\varphi}+e\sin E=\operatorname{arc.}\sin(\operatorname{tg}\frac{1}{2}\varphi^2\operatorname{tg}E)+e\sin E.$}. We shall pass over the expression for the maximum equation of centre as a series in powers of the eccentricity, which several authors have given. As an example, we append a table of the three maxima considered here for Juno, taking the eccentricity from the latest elements to be $=0{,}255\,4996$.
\begin{center}
\footnotesize
\setlength{\tabcolsep}{3pt}
\renewcommand{\arraystretch}{1.3}
\begin{tabular}{|c|c|c|c|c|}
\hline
Maximum & $E\quad\text{·}$ & $E-M$ & $v-E$ & $v-M$\\
\hline
$E-M$ & $90^0\ 0^{\prime}\ 0^{\prime\prime}$ & $14^0\,38^{\prime}\,20^{\prime\prime}{,}57$ & $14^0\,48^{\prime}\,11^{\prime\prime}{,}48$ & $29^0\,26^{\prime}\,32^{\prime\prime}{,}05$\\
$v-E$ & $82\ \ 32\ \ 9$ & $14\ \ 30\ \ 54{,}01$ & $14\ \ 55\ \ 41{,}79$ & $29\ \ 26\ \ 35{,}80$\\
$v-M$ & $86\ \ 14\ \ 40$ & $14\ \ 36\ \ 27{,}39$ & $14\ \ 53\ \ 49{,}57$ & $29\ \ 30\ \ 16{,}96$\\
\hline
\end{tabular}
\end{center}
